
\usepackage{tikz}
\usetikzlibrary{arrows,positioning,fit,shapes,decorations,decorations.pathmorphing,decorations.pathreplacing,calc,patterns,scopes,matrix,backgrounds}
\usepackage{booktabs}
%% tabularx lets a table size its own columns: l takes each column's
%% natural width and X absorbs the rest, so the table fills \textwidth
%% without any hand-computed widths to go stale.
\usepackage{tabularx}
\usepackage{color}
\usepackage[hyperindex,breaklinks]{hyperref}
\usepackage{helvet}
\usepackage{enumitem}
\usepackage[letterpaper,margin=1in]{geometry}
\usepackage{listings}
\usepackage{verbatim}
\usepackage[T1]{fontenc}

\setlength{\parindent}{0.0in}
\setlength{\parskip}{0.1in}

\newenvironment{apient}{\small\verbatim}{\endverbatim}

\newcommand{\apidesc}[1]{%
{\addtolength{\leftskip}{4em}%
#1\par\medskip}
}

\newcommand{\definedin}[1]{%
\textbf{Defined in:} \code{#1}
}

\newcommand{\code}[1]{{\small\texttt{#1}}}

%% Comments are meant to recede without becoming unreadable; a muted blue
%% of roughly the lightness the old lightgray had.
\definecolor{commentblue}{rgb}{0.40,0.55,0.75}

%% One style for every code listing in every manual.  Keep it here: an
%% \lstset outside a group is global and lasts to the end of the document,
%% so a stray one in an input file silently restyles every later listing.
%% A manual adds only its own emph list; everything else belongs here.
%%
%% breakatwhitespace keeps a wrap from landing inside an identifier, and
%% postbreak marks the continuation so a wrapped line cannot be mistaken
%% for one the author wrote that way.
\lstset{
  language={[GNU]C++},
  basicstyle=\ttfamily\small,
  keywordstyle=\ttfamily\color{purple},
  commentstyle=\ttfamily\color{commentblue},
  stringstyle=\ttfamily\color{orange},
  emphstyle=\color{blue},
  frame=leftline,
  numbers=left,
  numberstyle=\tiny,
  stepnumber=3,
  numbersep=5pt,
  showstringspaces=false,
  breaklines=true,
  breakatwhitespace=true,
  postbreak=\mbox{\textcolor{red}{$\hookrightarrow$}\space},
}
